\scanpage{81}
En poussant plus loin le calcul d'identification, on obtiendrait $c_3$ -- $c_n$ -- en fonction de $c_1,c_2$. Ce calcul est d'ailleurs simplifié considérablement en utilisant les formules suivantes, déduites de (4) par dérivation
\[
\boxed{\mathord{\unsure{\wp|p}}''=6\wp^2-\frac12g_2\qquad\wp'''=12\wp\wp'}\tag*{(6)}
\]
On obtient p. ex
\[
\boxed{c_3=\frac{c_1^2}3,\quad c_4=\frac3{11}c_1c_2}\tag*{(7)}
\]

\noindent\deleted{\unsure{2)}} 3) \fbox{\uline{Fonction $\zeta(u)$}}

\[
\zeta u-\frac1u=\int_0^u-\left(\wp z-\frac1{z^2}\right)dz\qquad\zeta'u=-\wp u\tag*{(7)}
\]
\[
\zeta u=\frac1u-\frac{c_2}{3}u^3\ldots-\frac{c_n}{2n+1}u^{2n+1}\ldots-\tag*{(8)}
\]
$\zeta$ est impaire, admet les périodes pour pôles simples, avec \deleted{\unsure{pôles}} \unsure{résidu|résidus} 1.
\[
\zeta z=\frac1z+\sum_{\omega\in{\unsure{\mathfrak P|\mathfrak B}}^{*}}\left(\frac1{z\unsure{-|+}\omega}+\frac1{\omega}+\frac{z}{\omega^2}\right)\tag*{(9)}
\]
$\zeta$ n'est pas elliptique (\unsure{voir plus loin}). --
\[
\boxed{\zeta(z+2\omega)-\zeta(z)=2\eta}\tag*{\textcircled{10}}
\]
($=\mathrm c^{\mathrm{te}}$ ; comme on voit en dérivant le premier membre).

Faisant $z=-\omega$, on obtient
\[
\eta=\zeta(\omega)\tag*{\textcircled{11}}
\]
; $\eta_1,\eta_2,\omega_1,\omega_2$ sont liés par une relation ; si $\omega_1,\omega_2$ sont en disposition directe, on a
\[
\eta_1\omega_2-\eta_2\omega_1=\frac{i\pi}{\unsure{2}}
\]

% archive sheet number: 80
